\section{Introduction}

Large language models have rapidly evolved from conversational assistants into autonomous software engineering agents capable of planning, implementing, debugging, and deploying complex infrastructure stacks~\citep{jimenez2024swebench,liu2024agentbench,frontiereng2026}.
Modern agentic development frameworks demonstrate increasingly sophisticated capabilities in repository-level code generation, multi-file reasoning, and tool orchestration.
Yet the evaluation methodologies for these systems remain rooted in a \emph{benchmark-driven paradigm}: isolated, static, short-horizon tasks where each evaluation unit is independent and the agent is reset after every attempt.

This paradigm has a fundamental limitation: \emph{real production engineering is not a collection of independent tasks}.
In practice, production infrastructure engineering involves long-horizon iterative development, multi-file and cross-module reasoning, dependency management, runtime debugging, tool orchestration, deployment validation, continuous integration pipelines, environment instability, resource constraints, and recovery from cascading failures.
Static benchmarks cannot capture these dynamics---they measure instantaneous coding ability, not engineering evolution.

\textbf{The gap.}
As a result, strong benchmark performance does not necessarily imply practical engineering utility.
Models optimized for benchmark-style coding may still fail in realistic production environments where runtime correctness, execution robustness, and adaptive recovery become critical.
We argue that the field needs to move \emph{beyond benchmarks} toward \emph{evaluation frameworks}---systems that assess agentic models in realistic, multi-stage production environments with controlled intermediate state injection, multi-dimensional metrics, and systematic failure diagnosis.

\textbf{Our approach.}
We present \eva{} (Evaluating Agentic Models in Production Systems), a production-grounded evaluation framework built upon the \yatcc{} integrated platform.
\eva{} is designed around four pillars:
\begin{enumerate}[nosep]
\item \textbf{Framework}: A unified evaluation architecture supporting heterogeneous LLM providers, agent SDKs (OpenHands, Claude Code, Codex CLI, Kimi CLI), and containerized execution backends, with standardized API access through \textsc{AIhub}.
\item \textbf{Workloads}: Realistic long-horizon engineering workloads spanning compiler construction---from lexical analysis through syntax parsing, IR generation, optimization, and backend code generation---organized as serial dependency chains where each task's output becomes the next task's input. Two difficulty levels (\yatcc{} and \yatcc{}-Hard) span a spectrum from scaffolded implementation to full from-scratch construction.
\item \textbf{Metrics}: A utility-oriented multi-dimensional evaluation system measuring end-to-end engineering quality (Mean Reward, Pass Score, Pipeline Score), deployment success, planning stability, debugging capability, execution cost (tokens, turns, retries, elapsed time), and runtime effectiveness through a five-category failure taxonomy.
\item \textbf{Resurrection}: A controlled intervention mechanism that injects verified intermediate artifacts after task failure, decomposing cascading failure into predecessor-unreachable and downstream-incapable, enabling node-level diagnostic evaluation even when the full pipeline is broken.
\end{enumerate}

Our contributions are:
\begin{enumerate}[nosep]
\item We propose \eva{} as a production-grounded evaluation \emph{framework} (not a benchmark) that moves beyond synthetic, isolated task evaluation toward persistent, runtime-observable, production-grounded assessment of agentic models.
\item We design realistic long-horizon engineering workloads with strict serial dependencies and verifiable intermediate states, spanning two difficulty levels that quantify the value of scaffolding in agentic coding tasks.
\item We introduce a unified evaluation architecture and utility-oriented metrics that support heterogeneous LLMs and agent frameworks, measuring not just outcome scores but process cost, failure patterns, and deployment-relevant efficiency.
\item We provide empirical evidence across 22 model configurations, four agent backends, and two workload difficulty levels, demonstrating that serial evaluation reveals failure patterns invisible to static benchmarks and that resurrection recovers substantial downstream diagnostic signal.
\end{enumerate}